\chapter{Finestra Operazioni sul Farm (FarmOperationsWindow)}

\section{Scopo}
Coordina il workflow operativo end-to-end verso il farm remoto: generazione input, upload, esecuzione batch e recupero risultati.

\section{Funzioni principali}
\begin{itemize}[leftmargin=1.5em]
  \item Inizializzazione ambiente remoto.
  \item Generazione e sincronizzazione file input.
  \item Lancio catene CALMET/CALPUFF/CALPOST.
  \item Post-processing e download output.
\end{itemize}

\section{Sottofinestre principali avviabili dall'utente}
\subsection{Dialog Permessi (open\_permissions\_dialog)}
Configura permessi operativi su file/cartelle remoti.

\subsection{MeteoWindow}
Apre la finestra meteo dedicata per parametri e preparazione dati meteorologici.

\subsection{ConfigPuntualeWindow}
Apre il dialog di configurazione puntuale tramite \texttt{ConfigPuntualeWindow.show\_dialog(...)}.

\subsection{Dialog Aggregazione}
Dialog per configurare aggregazione su risultati CSV.

\subsection{Dialog Post-Process}
Dialog per opzioni di post-processing e pipeline CALPOST.

\subsection{Dialog Timeseries}
Dialog per configurazione estrazione serie temporali (\texttt{\_show\_timeseries\_config\_dialog}).

\subsection{Dialog Download Risultati}
Dialog per selezionare e scaricare artefatti finali dal farm locale/remoto.

\section{File e cartelle coinvolte}
\begin{itemize}[leftmargin=1.5em]
  \item Legge i JSON in \texttt{temp\_config/}
  \item Usa template in \texttt{Working\_Files/scripts/}
  \item Genera input intermedi in \texttt{Outputs/}
  \item Popola \texttt{CALMET\_INP/}, \texttt{CALPUFF\_INP/}, \texttt{CALPOST\_INP/}
\end{itemize}
